\section{Transformer Optimizations}
\textit{Discuss the computational challenges of implementing Transformers in PyTorch. Suggest
optimization techniques such as gradient clipping, mixed precision training, and efficient
batching. Provide code examples for at least two techniques.}

While Transformer models offer strong performance across a
wide range of NLP tasks, they also introduce significant
computational challenges.
The self-attention mechanism has quadratic complexity with
respect to sequence length, leading to high memory usage
and increased training time.
In addition, deep Transformer architectures are prone to
training instabilities such as exploding gradients.
As a result, practical implementations in PyTorch often
require optimization techniques to improve stability,
efficiency, and scalability.

One common challenge when training Transformers is the
presence of large gradient values, particularly in deep
models or when using high learning rates.
Exploding gradients can cause unstable updates and prevent
convergence.
Gradient clipping is a simple yet effective technique to
address this issue by constraining the norm of the gradients
during backpropagation.
By enforcing an upper bound on the gradient norm, the model
is prevented from making excessively large parameter updates,
resulting in more stable training.
The use of gradient clipping in the Transformer training
loop is shown in \Cref{lst:transformer_optimizations}.

Another important optimization technique is mixed precision
training.
Transformers are computationally intensive and memory-hungry,
making them well-suited for half-precision arithmetic.
Mixed precision training combines 16-bit floating point
operations with 32-bit precision where necessary, reducing
memory consumption and improving throughput on modern GPUs.
To avoid numerical instability, loss scaling is applied so
that small gradient values are not lost due to reduced
precision.
PyTorch provides native support for mixed precision training
through automatic casting and gradient scaling, as
demonstrated in \Cref{lst:transformer_optimizations}.

In addition to these techniques, efficient batching plays a
crucial role in optimizing Transformer performance.
Grouping sequences of similar lengths within a batch reduces
the amount of padding required, which in turn lowers
unnecessary computation in the attention mechanism.
This can lead to substantial speedups and better utilization
of GPU memory, especially for large datasets with high
variance in sequence length.

The relevant portions of code implementing gradient clipping
and mixed precision training are provided in
\Cref{lst:transformer_optimizations}.
These optimizations are essential for training Transformer
models reliably and efficiently, particularly when scaling
to larger datasets and deeper architectures.
\begin{listing}
    \begin{minted}{Python}
scaler = torch.cuda.amp.GradScaler()

for src, labels, mask in dataloader:
    optimizer.zero_grad()

    with torch.cuda.amp.autocast():
        logits = model(src, mask)
        loss = criterion(logits, labels)

    scaler.scale(loss).backward()

    torch.nn.utils.clip_grad_norm_(model.parameters(), max_norm=1.0)

    scaler.step(optimizer)
    scaler.update()
    \end{minted}
    \caption{Optimizing Transformer Training with Gradient Clipping and Mixed Precision}
    \label{lst:transformer_optimizations}
\end{listing}

